\documentclass[12pt]{article}
\usepackage[utf8]{inputenc}
\usepackage[T1]{fontenc}
\usepackage{amsmath, amssymb, geometry, booktabs}
\usepackage{enumitem}
\geometry{margin=1in}
\setlength{\parindent}{0pt}
\renewcommand{\arraystretch}{1.2}

\title{Solutions -- Practice Problem Set 2}
\author{ECON 308 -- International Economic Relations}
\date{Fall 2025}

\begin{document}
\maketitle

\section*{Short Conceptual}

\textbf{1. H--O: labour endowment increase.} A higher $L$ (holding $K$ fixed) rotates the PPF outward toward the labour-intensive good. At given goods prices (small-country), by \textit{Rybczynski}, output of the labour-intensive good rises and the capital-intensive good falls. With endogenous prices (large-country), the relative supply of the labour-intensive good rises, tending to reduce its relative price.

\medskip
\textbf{2. Rybczynski Theorem.} Holding goods prices constant, an increase in one factor endowment increases output of the good using that factor intensively and reduces output of the other good (in a two-good, two-factor competitive economy with full employment).

\medskip
\textbf{3. Stolper--Samuelson in a labour-scarce country.} If Country A is capital-scarce (so $L$ relatively abundant), opening to trade raises the relative price of the labour-intensive good it exports. By Stolper--Samuelson, the real wage (return to the abundant factor) rises and the real rental rate (return to the scarce factor) falls.

\medskip
\textbf{4. Export- vs. import-biased growth.} Export-biased growth expands the supply of the export good, tends to lower its world relative price, and \emph{worsens} the home terms of trade (TOT). Import-biased growth reduces demand for imports (or raises supply of import-competing good), tends to raise the price of exports relative to imports, \emph{improving} the TOT.

\medskip
\textbf{5. TOT improvement in the Standard Trade Model.} An increase in $P_{\text{exports}}/P_{\text{imports}}$ pivots the trade line outward at the production point, allowing a higher indifference curve---so welfare rises in a distortion-free competitive economy.

\medskip
\textbf{6. External vs. internal economies.} \emph{External} economies: unit costs fall with \emph{industry} output (many firms; imperfect competition not required), generating potential multiple equilibria and geographic concentration. \emph{Internal} economies: unit costs fall with \emph{firm} output, implying imperfect competition (e.g., monopolistic competition).

\medskip
\textbf{7. Path dependence.} With external economies, an early, possibly accidental, concentration of production lowers unit costs, attracting more firms, further lowering costs. History (``accidents'') can lock in an industrial location even if fundamentals are symmetric.

\newpage

\section*{Analytical and Numerical Problems}

\subsection*{Problem 1: Long-Run Factor Adjustment (H--O)}

\textbf{Data (Home):} $K_H=400,\ L_H=800$. Sectoral factor intensities: textiles (labour-intensive) $K_T/L_T=0.25$, steel (capital-intensive) $K_S/L_S=2$.

\begin{enumerate}[label=(\alph*)]
\item \textbf{Relative factor abundance.} Home's aggregate ratio is
\[
\left(\frac{K}{L}\right)_H=\frac{400}{800}=0.5.
\]
\emph{Interpretation:} By itself, $0.5$ does not establish abundance; abundance is \emph{relative across countries}. If Foreign has $(K/L)_F>0.5$, then Home is \textbf{labour-abundant}; if $(K/L)_F<0.5$, Home is \textbf{capital-abundant}. In what follows, we adopt the standard case that Home is labour-abundant (i.e., $(K/L)_F>0.5$).

\item \textbf{Export pattern.} By the H--O theorem, a labour-abundant country exports the labour-intensive good $\Rightarrow$ \textbf{Home exports textiles}.

\item \textbf{Factor price effects (Stolper--Samuelson).} Opening to trade raises the relative price of the labour-intensive good that Home exports. Thus the \emph{real wage} rises and the \emph{real rental rate} falls (measured in terms of both goods).

\item \textbf{10\% rise in $P_T$ (world price of textiles).} By Stolper--Samuelson magnification, $w$ rises more than proportionately relative to the numéraire good used intensively in textiles, while $r$ falls in real terms. Labour moves toward textiles in the new long-run equilibrium, and outputs reallocate accordingly.
\end{enumerate}

\bigskip

\subsection*{Problem 2: Standard Trade Model and Welfare}

\textbf{PPF:} $F^2+2E^2=200$; \textbf{world relative price:} $P_E/P_F=2$.

\medskip
\textbf{(a) Value-maximizing production.} Maximize $V=P_F F+P_E E$. Normalize $P_F=1$, so $V=F+2E$ subject to $F^2+2E^2=200$.

Lagrangian:
\[
\mathcal{L}=F+2E+\lambda\,(200-F^2-2E^2).
\]
First-order conditions (FOCs):
\[
\frac{\partial \mathcal{L}}{\partial F}=1-2\lambda F=0 \ \Rightarrow\ F=\frac{1}{2\lambda},\qquad
\frac{\partial \mathcal{L}}{\partial E}=2-4\lambda E=0 \ \Rightarrow\ E=\frac{1}{2\lambda}.
\]
Hence $F=E$ at the optimum. Substitute into the PPF:
\[
F^2+2F^2=3F^2=200 \ \Rightarrow\ F=E=\sqrt{200/3}\approx 8.165.
\]
\[
\boxed{(F_P,E_P)=\left(\sqrt{\tfrac{200}{3}},\ \sqrt{\tfrac{200}{3}}\right)\approx(8.165,\,8.165)}.
\]
Production value: $V=F+2E=3F\approx 24.495$.

\medskip
\textbf{(b) Consumption with $F=E$ preferences and trade pattern.} Consumption must satisfy the budget line $F+2E=V$ at world prices and lie on the preference ray $E=F$. Thus $F+2F=3F=V \Rightarrow F_C=V/3=\sqrt{200/3}$ and $E_C=F_C$. Therefore,
\[
\boxed{(F_C,E_C)=(F_P,E_P)}.
\]
\emph{Trade pattern:} Zero trade volumes (the country happens to produce at its most-preferred bundle at these prices).

\medskip
\textbf{(c) Gains from trade vs. autarky.} Because $(F_P,E_P)$ already lies at the tangency of the world price line with the PPF \emph{and} preferences pick $F=E$, autarky and free-trade consumption coincide. Hence there are no additional gains-from-trade at this price vector and preference ray (a knife-edge case). In general, with different preferences, trade would move consumption to a higher indifference curve.

\medskip
\textbf{(d) TOT improvement (rise in $P_E/P_F$).} If electronics is the export good and $P_E/P_F$ rises, the world price line becomes steeper, rotating outward at the production point. Firms re-optimize (more $E$, less $F$), and the budget line through the new production point supports a \emph{higher} indifference curve. Thus welfare \emph{increases}. (A standard STM diagram shows the new, steeper price line tangent to the PPF at a point with higher $E/F$, and a higher indifference curve.)

\bigskip

\subsection*{Problem 3: External Economies and Location }

\textbf{Costs and demand.} Unit cost in a country with industry output $Q$:
\[
C(Q)=20-0.05Q, \quad \text{(external economies: costs fall with industry output)}.
\]
World demand: $Q^W=400-5P \ \Rightarrow\ P=80-0.2Q^W$.

\medskip
\textbf{(a) Price--quantity relationship.} From demand:
\[
\boxed{P=80-0.2Q^W}.
\]

\medskip
\textbf{(b) Symmetric production (each makes half).} If A and B each produce $Q^W/2$, their unit costs are
\[
C_A=C_B=20-0.05\left(\frac{Q^W}{2}\right)=20-0.025Q^W.
\]
Under competition, $P$ equals the common unit cost. Equate with inverse demand:
\[
80-0.2Q^W = 20-0.025Q^W 
\ \Rightarrow\ 60=0.175\,Q^W \ \Rightarrow\ Q^W=\frac{60}{0.175}\approx 342.857.
\]
Then
\[
P=80-0.2(342.857)\approx 11.4286, \quad C=20-0.025(342.857)\approx 11.4286.
\]
\[
\boxed{Q^W\approx 342.86,\quad P\approx 11.43.}
\]

\medskip
\textbf{(c) Positive feedback and specialization.} Suppose A produces slightly more: $Q_A=\frac{Q^W}{2}+\varepsilon$, $Q_B=\frac{Q^W}{2}-\varepsilon$.
\[
C_A=20-0.05\!\left(\tfrac{Q^W}{2}+\varepsilon\right)=20-0.025Q^W-0.05\varepsilon,
\]
\[
C_B=20-0.05\!\left(\tfrac{Q^W}{2}-\varepsilon\right)=20-0.025Q^W+0.05\varepsilon.
\]
Hence
\[
C_A - C_B = -0.1\,\varepsilon < 0 \quad (\varepsilon>0).
\]
A's unit cost falls \emph{below} B's by $0.1\varepsilon$, letting A undercut B, capture more market share, raise $Q_A$ further, and lower costs again—\textbf{circular causation} that drives the industry toward \emph{complete specialization} in one country.

\medskip
\textbf{(d) Trade between identical countries.} Even if A and B are identical ex ante, historical accidents (small initial $\varepsilon$) can select a single production location. The specialized country exports to the other at the common world price; thus identical countries \emph{do} trade due to external economies and path dependence (not factor endowments).

\medskip
\textbf{(e) Policy and permanent relocation.} A temporary production subsidy (or infant-industry protection) in A raises $Q_A$, lowers $C_A$, and can tip the industry into A. After scale builds, the subsidy can be removed; the new low-cost location persists (hysteresis) because exiting the cluster would raise costs.


\end{document}
